% GENERATED FILE. Edit canonical lessons, then run npm run generate:books.
% canonical-source-hash: fnv1a64:75cd6c11acdee771
% canonical-lessons: FR-C511-flute, FR-C511-violoncelle, FR-C511-trompette, FR-C511-accordeon, FR-C511-accorder

\chapter{Instruments de musique}
\label{ch:fr-instruments-de-musique}
\hlchaptermodality{511}

\begin{chapteropening}
\textbf{By the end of this chapter:} \emph{I can talk about the instrument I play.}

Complete the last lesson of chapter 511: I can talk about the instrument I play.
\end{chapteropening}

\section[la flûte]{la flûte — flute; recorder}

\label{lesson:FR-C511-flute}



\begin{quote}
Before the new one: say the French for ice rink, then the French for slope.
\end{quote}

\subsection*{You'll want to know: la flûte}
\textbf{la flûte} — \textquotedblleft{}flute; recorder\textquotedblright{}.

\begin{quote}
Au collège, tous les élèves apprennent la flûte.
\end{quote}

\begin{quote}
\emph{At secondary school all the pupils learn the recorder.}
\end{quote}

\emph{La flûte à bec} — the recorder; \emph{la flûte traversière} — the flute. Also a tall champagne glass.

\begin{cousinweb}
\emph{la flûte} goes back to Old French \emph{flaüte} or \emph{fleute}, probably from Old Occitan \emph{flaüt} (EtymDB). The origin is not settled; it may imitate the sound.
\end{cousinweb}

\begin{grammarlens}[title={What you've built}]
One more word for telling, explaining and describing at B1.
\end{grammarlens}

\subsection*{Your turn}
\begin{itemize}
\raggedright
  \item \emph{Say it:} \emph{la flûte}
  \item \emph{Say it:} \emph{la flûte}, once more
  \item \emph{Say it:} a sentence of your own with \emph{la flûte}
  \item \emph{Recall:} say the French for ice rink, then the French for slope, then say \emph{la flûte} again
\end{itemize}

\begin{tcolorbox}[breakable,colback=teal!4,colframe=teal!35!black,title={Before you move on}]
What does la flûte mean? (\textquotedblleft{}Flute; recorder\textquotedblright{}.) Say it once more.
\end{tcolorbox}
\section[le violoncelle]{le violoncelle — cello}

\label{lesson:FR-C511-violoncelle}



\begin{quote}
Before the new one: say the French for slope, then the French for flute.
\end{quote}

\subsection*{You'll want to know: le violoncelle}
\textbf{le violoncelle} — \textquotedblleft{}cello\textquotedblright{}.

\begin{quote}
Ma sœur joue du violoncelle depuis l'âge de huit ans.
\end{quote}

\begin{quote}
\emph{My sister has played the cello since she was eight.}
\end{quote}

\emph{Jouer du violoncelle}. The player is \emph{le} or \emph{la violoncelliste}.

\begin{cousinweb}
\emph{le violoncelle} was borrowed from Italian \emph{violoncello}, a small form of \emph{violone} \textquotedblleft{}a large viol\textquotedblright{}, itself built on \emph{viola}.
\end{cousinweb}

\begin{grammarlens}[title={What you've built}]
One more word for telling, explaining and describing at B1.
\end{grammarlens}

\subsection*{Your turn}
\begin{itemize}
\raggedright
  \item \emph{Say it:} \emph{le violoncelle}
  \item \emph{Say it:} \emph{le violoncelle}, once more
  \item \emph{Say it:} a sentence of your own with \emph{le violoncelle}
  \item \emph{Recall:} say the French for slope, then the French for flute, then say \emph{le violoncelle} again
\end{itemize}

\begin{tcolorbox}[breakable,colback=teal!4,colframe=teal!35!black,title={Before you move on}]
What does le violoncelle mean? (\textquotedblleft{}Cello\textquotedblright{}.) Say it once more.
\end{tcolorbox}
\section[la trompette]{la trompette — trumpet}

\label{lesson:FR-C511-trompette}



\begin{quote}
Before the new one: say the French for flute, then the French for cello.
\end{quote}

\subsection*{You'll want to know: la trompette}
\textbf{la trompette} — \textquotedblleft{}trumpet\textquotedblright{}.

\begin{quote}
Mon voisin joue de la trompette le dimanche matin.
\end{quote}

\begin{quote}
\emph{My neighbour plays the trumpet on Sunday mornings.}
\end{quote}

\emph{Le} or \emph{la trompettiste} — the trumpeter.

\begin{cousinweb}
\emph{la trompette} is built on \emph{trompe} \textquotedblleft{}horn\textquotedblright{} with the small ending \emph{-ette}; \emph{trompe} comes from a Frankish word \emph{trumpa}.
\end{cousinweb}

\begin{grammarlens}[title={What you've built}]
One more word for telling, explaining and describing at B1.
\end{grammarlens}

\subsection*{Your turn}
\begin{itemize}
\raggedright
  \item \emph{Say it:} \emph{la trompette}
  \item \emph{Say it:} \emph{la trompette}, once more
  \item \emph{Say it:} a sentence of your own with \emph{la trompette}
  \item \emph{Recall:} say the French for flute, then the French for cello, then say \emph{la trompette} again
\end{itemize}

\begin{tcolorbox}[breakable,colback=teal!4,colframe=teal!35!black,title={Before you move on}]
What does la trompette mean? (\textquotedblleft{}Trumpet\textquotedblright{}.) Say it once more.
\end{tcolorbox}
\section[l'accordéon]{l'accordéon — accordion}

\label{lesson:FR-C511-accordeon}



\begin{quote}
Before the new one: say the French for cello, then the French for trumpet.
\end{quote}

\subsection*{You'll want to know: l'accordéon}
\textbf{l'accordéon} — \textquotedblleft{}accordion\textquotedblright{}.

\begin{quote}
Au bal du village, un vieil homme jouait de l'accordéon.
\end{quote}

\begin{quote}
\emph{At the village dance an old man was playing the accordion.}
\end{quote}

Masculine. A typical sound of the French \emph{bal musette}.

\begin{cousinweb}
\emph{l'accordéon} was borrowed from German \emph{Akkordion} (EtymDB), built on \emph{Akkord} \textquotedblleft{}chord\textquotedblright{}, which itself was borrowed from French \emph{accord}; \emph{accord} goes back to Latin \emph{cor} \textquotedblleft{}heart\textquotedblright{} (Lewis \& Short).
\end{cousinweb}

\begin{grammarlens}[title={What you've built}]
One more word for telling, explaining and describing at B1.
\end{grammarlens}

\subsection*{Your turn}
\begin{itemize}
\raggedright
  \item \emph{Say it:} \emph{l'accordéon}
  \item \emph{Say it:} \emph{l'accordéon}, once more
  \item \emph{Say it:} a sentence of your own with \emph{l'accordéon}
  \item \emph{Recall:} say the French for cello, then the French for trumpet, then say \emph{l'accordéon} again
\end{itemize}

\begin{tcolorbox}[breakable,colback=teal!4,colframe=teal!35!black,title={Before you move on}]
What does l'accordéon mean? (\textquotedblleft{}Accordion\textquotedblright{}.) Say it once more.
\end{tcolorbox}
\section[accorder]{accorder — to tune (an instrument); to grant}

\label{lesson:FR-C511-accorder}



\begin{quote}
Before the new one: say the French for trumpet, then the French for accordion.
\end{quote}

\subsection*{You'll want to know: accorder}
\textbf{accorder} — \textquotedblleft{}to tune (an instrument); to grant\textquotedblright{}.

\begin{quote}
Avant le concert, le guitariste accorde sa guitare.
\end{quote}

\begin{quote}
\emph{Before the concert the guitarist tunes his guitar.}
\end{quote}

Also to grant: \emph{accorder une pause}. \emph{S'accorder} — to agree.

\begin{cousinweb}
\emph{accorder} is inherited from Vulgar Latin \emph{accordāre} \textquotedblleft{}to bring into agreement\textquotedblright{} (EtymDB), built on Latin \emph{cor} \textquotedblleft{}heart\textquotedblright{} (Lewis \& Short): to make hearts agree.
\end{cousinweb}

\begin{grammarlens}[title={What you've built}]
One more word for telling, explaining and describing at B1.
\end{grammarlens}

\subsection*{Your turn}
\begin{itemize}
\raggedright
  \item \emph{Say it:} \emph{accorder}
  \item \emph{Say it:} \emph{accorder}, once more
  \item \emph{Say it:} a sentence of your own with \emph{accorder}
  \item \emph{Recall:} say the French for trumpet, then the French for accordion, then say \emph{accorder} again
\end{itemize}

\begin{tcolorbox}[breakable,colback=teal!4,colframe=teal!35!black,title={Before you move on}]
What does accorder mean? (\textquotedblleft{}To tune (an instrument); to grant\textquotedblright{}.) Say it once more.
\end{tcolorbox}
